\color{unidarkgreen}{\subsection{Технологические защиты}}
\color{black}

\begin{enumerate}

\item
Технологические защиты от повышения давления, температуры масла и обмоток предназначены для предотвращения повреждения оборудования вследствие нарушения режима его работы, обеспечивая надёжную и безопасную работу ШР. Для ШР предусматриваются функциональные блоки приема и обработки сигналов от датчиков температуры масла и обмотки.

\item
Предусмотрена возможность приема сигналов срабатывания сигнального контакта ДТм ШР (ДТо ШР), отключающего контакта ДТм ШР  (ДТо ШР) и низкой изоляции цепи отключающего контакта ДТм ШР (ДТо ШР), по разным цепям от двух источников ПДС1 и ПДС2 пофазно. Для приема сигналов от ПДС1 предусмотрены входные сигналы \textit{«Ср. сигн контакта ДТм ф. «А» ШР (ДТо ШР ф. «А») 1 цепь»}, \textit{«Ср. авар контакта ДТм ф. «А» ШР (ДТо ШР ф. «А») 1 цепь»} и \textit{«Низкая изол. ДТм ф. «А» ШР  (ДТо ШР ф. «А») авар 1 цепь»}, а для приема сигналов от ПДС2 предусмотрены входные сигналы \textit{«Ср. сигн контакта ДТм ШР ф. «А» (ДТо ШР ф. «А») 2 цепь»}, \textit{«Ср. авар контакта ДТм ШР ф. «А» (ДТо ШР ф. «А») 2 цепь»} и \textit{«Низкая изол. ДТм ШР ф. «А» (ДТо ШР ф. «А» ) авар 2 цепь»}. Если сигналы срабатывания ДТм ШР ф. «А» (ДТо ШР м) заводятся непосредственно со шкафа ШОМ ШР конфигурируются только входные сигналы \textit{«Ср. сигн контакта ДТм ШР ф. «А» (ДТо ШР ф. «А») 1 цепь»}, \textit{«Ср. авар контакта ДТм ШР ф. «А» (ДТо ШР ф. «А») 1 цепь»} и \textit{«Низкая изол. ДТм ШР ф. «А» (ДТо ШР ф. «А») авар 1 цепь»}, а действие ДТм ШР ф. «А» (ДТо ШР ф. «А») по второй цепи сигнального и аварийного контакта выводится из работы с помощью уставок \textbf{«ДТм ШР (ДТо ШР) сигн 2 цепь»} и \textbf{«ДТм ШР (ДТо ШР) авар 2 цепь»}.

\item
Ввод в работу ДТм ШР (ДТо ШР) осуществляется с помощью уставок \textbf{«ДТм ШР (ДТо ШР) сигн 1 цепь»}, \textbf{«ДТм ШР (ДТо ШР) сигн 2 цепь»}, \textbf{«ДТм ШР (ДТо ШР) авар 1 цепь»} и \textbf{«ДТм ШР (ДТо ШР) авар 2 цепь»}. Предусмотрена возможность действия срабатывания аварийного контакта ДТм ШР (ДТо ШР) на отключение трансформатора с помощью уставки \textbf{«Откл от ДТм ШР (ДТо ШР) авар»}. Для фаз «B»
и «С» аналогично. 

\item
Оперативный ввод/вывод ТЗ производится с помощью БУ (блока управления функциями) \textit{«Вывод ТЗ Т»}.

\item
Непрерывный контроль изоляции цепи аварийного контакта ДТм (ДТо) осуществляется с помощью реле контроля изоляции (РКИ). В случае снижения сопротивления изоляции РКИ подает сигнал \textit{«Низкая изол. ДТм ШРм ф. «А» (ДТо ШР ф. «А») авар 1 цепь»} (\textit{«Низкая изол. ДТм ШР ф. «А» (ДТо ШР ф. «А») авар 2 цепь»}) для блокировки срабатывания ДТм ШР ф. «А» (ДТо ШР ф. «А») от аварийного контакта. Время срабатывания блокировки от сигналов РКИ задается с помощью уставки \textbf{«Тбл ТЗ ШР»}. Ввод блокировки ДТм ШР (ДТо ШР) от РКИ осуществляется с помощью уставки \textbf{«Блок ТЗ ШР откл»}. Для сброса блокировки от РКИ, используется команда управления \textit{«Возврат ТЗ ШР»}.

\item 
Уставки ТЗ ШР приведены в таблице \ref{app_set:tz}.

\item 
Алгоритм работы сигнальной ступени ДТм (ДТо) фазы А приведен на рисунке \ref{pict_tz:TZ}, алгоритмы работы сигнальных ступеней ДТм(ДТо) других фаз выполнены аналогично.


\medskip
\begin{figure}[H]
\centering
\includegraphics[width=1\textwidth,height=1\textheight,keepaspectratio]{\fbpath/11.01 ШР/161.1061 ТЗ/img/TZ_sign.pdf}
\captionof{figure}{Функциональная схема логики сигнальной ступени ДТм фазы А}\label{pict_tz:TZ}
\end{figure}

\item 
Алгоритм работы отключающей ступени ДТм (ДТо) фазы А приведен на рисунке \ref{pict_tz:TZotkl}, алгоритмы работы отключающих ступеней ДТм(ДТо) других фаз выполнены аналогично.

\medskip
\begin{figure}[H]
\centering
\includegraphics[width=1\textwidth,height=1\textheight,keepaspectratio]{\fbpath/11.01 ШР/161.1061 ТЗ/img/TZ_otkl.pdf}
\captionof{figure}{Функциональная схема логики отключающей ступени ДТм фазы А}\label{pict_tz:TZotkl}
\end{figure}


\end{enumerate}